% =====================================================================
%  PROTOCOLE TRIADE v0.5.1 — Manuscrit de spécification
%  Corps du document : Tectonic via pdf.py convert.latex.
%  Couverture générée séparément (cover.html, Template 03) et fusionnée
%  en amont — AUCUNE page de titre ici. Le sommaire est la première page
%  du corps.
%  Chaîne née spec_v040 → spec_v050 (D25/D44) → spec_v051 (D75).
% =====================================================================
\documentclass[11pt,a4paper]{article}

% ---------- Polices ----------
\usepackage{fontspec}
\setmainfont{Liberation Serif}
\setsansfont{Liberation Sans}
\setmonofont{Liberation Mono}
\usepackage[french]{babel}
\addto\captionsfrench{\renewcommand{\tablename}{Tableau}}

% ---------- Mathématiques ----------
\usepackage{amsmath,amssymb,amsthm}

% ---------- Mise en page ----------
\usepackage{geometry}
\geometry{a4paper, top=2.4cm, bottom=2.4cm, left=2.2cm, right=2.2cm, headheight=14pt}
\usepackage{graphicx}
\usepackage{xcolor}
\definecolor{liaison}{HTML}{3681A6}
\definecolor{grisdoux}{HTML}{555F66}
\usepackage{booktabs}
\usepackage{tabularx}
\newcolumntype{Y}{>{\raggedright\arraybackslash}X}
\usepackage{multirow}
\usepackage{adjustbox}
\usepackage{enumitem}
\setlist{noitemsep, topsep=2pt, leftmargin=1.4em}
\usepackage{microtype}

% ---------- Théorèmes ----------
\theoremstyle{definition}
\newtheorem{regle}{Règle}

% ---------- Listings ----------
\usepackage{listings}
\lstset{
  basicstyle=\footnotesize\ttfamily,
  keywordstyle=\bfseries,
  commentstyle=\itshape\color{grisdoux},
  breaklines=true,
  showstringspaces=false,
  columns=fullflexible,
  keepspaces=true,
  frame=single,
  framerule=0.3pt,
  rulecolor=\color{black!25},
  xleftmargin=4pt, xrightmargin=4pt,
  aboveskip=6pt, belowskip=6pt,
}

% ---------- Titres ----------
\usepackage{titlesec}
\titleformat{\section}
  {\Large\sffamily\bfseries\color{liaison}}{\thesection}{0.6em}{}
\titleformat{\subsection}
  {\large\sffamily\bfseries}{\thesubsection}{0.6em}{}
\titlespacing*{\section}{0pt}{14pt plus 4pt minus 2pt}{7pt plus 2pt minus 2pt}

% ---------- En-têtes / pieds ----------
\usepackage{fancyhdr}
\pagestyle{fancy}
\fancyhf{}
\fancyhead[L]{\small\itshape Protocole TRIADE — spécification v0.5.1}
\fancyfoot[C]{\small\thepage}
\renewcommand{\headrulewidth}{0.3pt}

% ---------- Divers ----------
\raggedbottom
\widowpenalty=10000
\clubpenalty=10000
\setcounter{tocdepth}{2}

% hyperref : TOUJOURS en dernier
\usepackage[colorlinks=true, linkcolor=liaison, citecolor=grisdoux,
            urlcolor=liaison, bookmarks=true, bookmarksnumbered=true,
            unicode=true]{hyperref}

\hypersetup{
  pdftitle={Protocole TRIADE v0.5.1 — spécification technique},
  pdfauthor={Alexis FR-LOPEZ ; Agent Z.ai (GLM)},
  pdfsubject={Gouvernance vérifiable du développement logiciel assisté par IA},
  pdfkeywords={TRIADE, agent-first, contrat, vérification, gouvernance, profils}
}

\newcommand{\code}[1]{\texttt{#1}}

\begin{document}

% Le sommaire est la première page du corps (la couverture est fusionnée
% en amont). Conformité avec la règle « TOC = première page ».
\tableofcontents
\newpage

% =====================================================================
%  CH 1-4 : Vision · Couche I et taxonomie · Couches et traçabilité · Carte
% =====================================================================
\section{Vision et objectifs}\label{sec:vision}

Le développement logiciel est entré dans l’ère \emph{agent-first} : une part croissante du
code est écrite, relue et modifiée par des modèles de langage plutôt que par des humains.
Cette bascule remet en question une hypothèse centrale de nos outils — que le code source
est optimisé pour la lecture humaine. TRIADE est la réponse disciplinée à ce réexamen :
plutôt que de supprimer les langages (une impossibilité de principe, puisqu’un modèle de
langage manipule des symboles textuels et que tout format binaire rendu lisible par un
agent constitue de facto un langage), TRIADE déplace le langage vers une fine couche de
spécification maintenue par l’agent, tandis que l’artefact livré est du code machine natif
pur. Le compilateur disparaît du pipeline de livraison : ce que l’agent écrit est ce qui
s’exécute — et rien d’autre.

Le nom encode l’architecture : trois couches indissociables — Intention (I), Binaire (B),
Carte (C) — unies par un verrou croisé cryptographique. La couche I contient le quoi et le
pourquoi : contrats typés, pré et post-conditions, invariants, cas d’échec, vecteurs de
vérification. La couche B contient le comment : des octets scellés, versionnés, à adresses
stables. La couche C est l’index naviguable qui lie les deux et sert d’unique interface de
développement à l’agent.

La v0.4.0 tire les enseignements de la première validation à l’échelle — le recodage
complet de RoomEq (démonstration v0.3) — et les convertit en normes : ce que la
démonstration a pratiqué (reverse-specs, pré-versions, partage de fixtures, mesure avant
optimisation) devient exigence ; ce que son analyse des défauts a révélé devient
taxonomie ; ce que son étude de cas établit et ne prouve pas reste explicitement hors
périmètre de preuve.

La v0.5.0 est la révision du retour du terrain. Les jalons M0 (chaîne artefact), M1
(mini-libc scellée) et M2 (patch à chaud en production, avec retour arrière) sont
franchis — chacun sur validation terrain Windows 11, chacun scellé par une balise
finale immuable (0.2.0, 0.3.0). La campagne M2 a coûté huit classes de défauts
invisibles à l’émulation (§\ref{sec:taxonomieX}) — le prix d’une leçon : un oracle
valide un artefact contre un modèle d’environnement, et la complétude du modèle ne se
démontre qu’en exécutant le binaire sur le système réel. Ce que le terrain a appris
devient norme : la boucle d’évidence (§\ref{sec:terrain}), la taxonomie X,
l’instrumentation contractualisée — et l’ouverture du jalon M3, contrats renforcés
sur entrées non fiables.

\subsection{Fonction objectif hiérarchisée — invariant de révision}\label{sec:objectif}

\begin{center}
\large\bfseries Sécurité $>$ Performance $>$ Optimisation $>$ Coût en tokens
\end{center}

\medskip

Arbitrage lexicographique : un objectif n’est jamais sacrifié au profit d’un objectif de
rang inférieur. La v0.4.0 consigne cette hiérarchie comme \textbf{invariant de révision}
(D19) : aucune révision future du protocole ne peut la modifier sans invalider
explicitement le protocole entier — pas l’ordre, pas les rangs, pas la non-compensation.

\begin{table}[htbp]
\centering\small
\caption{La fonction objectif et ses métriques de validation.}
\label{tab:objectif}
\begin{tabularx}{\textwidth}{@{}l Y Y@{}}
\toprule
\textbf{Objectif} & \textbf{Signification concrète} & \textbf{Métrique de validation} \\
\midrule
Sécurité & Sûreté mémoire validée par oracle ; intégrité carte $\leftrightarrow$ binaire ;
durcissement au chargement & 0 corruption détectée ; tout échec de verrou bloque la
session ; échec au lancement = arrêt propre \\
Performance & 100 \% code machine natif ; zéro abstraction intermédiaire &
Comparaison par bloc scellé face à une référence compilée \\
Optimisation & Descente à l’ISA sur hot paths identifiés — jamais sans mesure (R11) &
Gain mesuré et journalisé par bloc optimisé \\
Coût en tokens & Édition ciblée via la carte ; jamais de lecture intégrale du binaire &
Budgets par type d’opération (§\ref{sec:carte}) \\
\bottomrule
\end{tabularx}
\end{table}

\subsection{Périmètre et non-objectifs}\label{sec:perimetre}

Cinq non-objectifs. Premièrement, TRIADE ne prétend pas remplacer les langages dans
l’industrie : c’est une architecture de laboratoire. Deuxièmement, la furtivité face aux
EDR n’est pas recherchée : l’outil cible est un utilitaire de pentest pour tests autorisés
en laboratoire, appelant les API du système via les interfaces standards du profil.
Troisièmement, l’universalité est structurelle mais progressive : chaque profil doit être
instancié et validé (M6, M7) avant d’être déclaré supporté — le protocole ne promet pas la
portabilité, il la vérifie. Quatrièmement, aucune couche d’abstraction à l’exécution
(bytecode portable, machine virtuelle, traduction à la volée) n’est introduite : elle
trahirait la fonction objectif Performance. Cinquièmement, la couche de gouvernance
validée par la démonstration et le cœur natif du protocole (notation, verrou
cryptographique sur octets, patch à chaud) restent des niveaux de preuve distincts : les
jalons M0 à M7 demeurent la charge de validation du cœur natif.

\subsection{Ce que change la v0.4.0}\label{sec:nouveautes}

Dix-sept évolutions, toutes tracées dans la correspondance du §\ref{sec:correspondance} et
dans le journal des versions, toutes rattachées aux décisions D19 à D24 : normalisation de
la couche I, taxonomie des défauts, traçabilité I $\leftrightarrow$ C $\leftrightarrow$
B $\leftrightarrow$ tests, norme de fragmentation, localisation par intention, R10
renforcée, R11 nouvelle, suite de conformité par profil, parité inter-profils,
pré-versions, reverse-specs, journal d’optimisation, modèle de menace étendu, journal des
décisions normalisé, tableau de bord de conformité. Aucune règle existante n’est
affaiblie ; aucune couche n’est ajoutée à l’exécution ; la hiérarchie de la fonction
objectif est inchangée.

\subsection{Ce que change la v0.5.0}

Huit évolutions, toutes tracées dans la correspondance (§\ref{sec:evolutions050}) et le
journal des versions, toutes rattachées aux décisions D26 à D45 : la validation
terrain devient une couche de vérification normée — boucle d’évidence, journal
d’exécution, instrumentation contractualisée, diff chirurgical (§\ref{sec:terrain}) ;
la taxonomie des défauts s’étend des classes T (contrat non écrit) aux classes X —
défauts invisibles à l’émulation, huit classes nommées par la campagne M2
(§\ref{sec:taxonomieX}) ; l’état des jalons est constaté — M0, M1 et M2 franchis sur
validations terrain Windows 11, balises finales 0.2.0 et 0.3.0 posées
(§\ref{sec:jalons}) ; le budget de localisation est corrigé à 64 tokens (D31) et les
budgets affichent leurs valeurs mesurées ; les volumes d’oracle effectifs sont
consignés ; les preuves terrain sont archivées brutes avec leurs empreintes ; le
corpus de démonstration (huit itérations instrumentées) devient l’occurrence de
référence de la méthode ; le jalon M3 est ouvert — contrats renforcés sur entrées non
fiables et fuzzing à propriétés (D45). Aucune règle n’est modifiée ; aucune couche
n’est ajoutée à l’exécution ; la hiérarchie de la fonction objectif est inchangée.

\subsection{Ce que change la v0.5.1}\label{sec:nouveautes051}

Un composant nouveau et deux constats, tous tracés dans la correspondance
(§\ref{sec:evolutions051}) et le journal des versions, rattachés à la décision D75 : le
\textbf{coffre-fort des phrases maîtresses} — conteneur chiffré TRIADE-CF1 centralisant
toutes les phrases maîtresses par projet, sur la pile cryptographique de la carte au
repos (Argon2id aux planchers, HKDF à contexte dédié, XChaCha20-Poly1305), doté d’un
autotest de falsification et de règles d’export/import (continuité,
§\ref{sec:gouvernance} ; FORMAT\_CARTE § 11) ; l’état des jalons est constaté — M3
franchi le 28 octobre 2026 sur la campagne million (D72), la conformité permanente
(D73) et le scellage de la balise finale 0.4.0 (D74) ; le modèle de menace gagne la
ligne du vol du coffre-fort (§\ref{sec:menaces}). Aucune règle n’est modifiée ; aucune
couche n’est ajoutée à l’exécution ; la hiérarchie de la fonction objectif est inchangée.

\section{La couche I : format normalisé des spécifications}\label{sec:coucheI}

La v0.3 décrivait la couche I comme « contrats légers » (D3), sans en fixer le gabarit —
la légèreté demeure, la structure devient normative (D19). Une spécification I n’est pas un
document libre : c’est un artefact diffable dont la structure conditionne la vérification.

\subsection{Identifiants stables}

Chaque spécification reçoit à sa création un identifiant \code{<DOM>-<NN>} : un domaine
parmi une liste figée par projet et un numéro séquentiel. L’identifiant est immuable et
jamais réutilisé ; un module supprimé laisse son identifiant vacant à jamais. C’est
l’équivalent exact, côté intention, de l’adresse stable côté binaire (R7) : les références
croisées — dépendances, carte, tests — ne cassent jamais. Le couple (identifiant,
révision) est la clé de traçabilité du §\ref{sec:traceabilite}.

\subsection{Gabarit obligatoire}

\begin{lstlisting}
<DOM>-<NN> . <nom court> . revision <R>
Intention : <une phrase — le pourquoi>
Contrat :
  pre  : P1..Pn    (preconditions numerotees, une phrase chacune)
  post : Q1..Qn    (postconditions numerotees)
  inv  : I1..In    (invariants)
Echecs :
  E1 : <condition> -> <sanction nommee : exception | arret propre | repli>
Vecteurs :
  V1 : <vecteur dore | propriete differentielle | test de rejet>
Fixtures    : <references aux fixtures contractuelles partagees>
Provenance  : ecrite-avant-code | reverse-spec de <binaire scelle @ tag>
Decisions   : <Dxx...>
\end{lstlisting}

Sept blocs, tous obligatoires — les échecs et vecteurs ne peuvent être vides que pour un
module sans effet observable, situation rare à justifier dans le champ Intention. Trois
règles de rédaction conditionnent la diffabilité exigée par R2 : \textbf{une exigence par
phrase numérotée} (le diff montre exactement quelle exigence change) ; \textbf{aucune
connaissance dupliquée} (renvoi par identifiant, sinon défaut T5) ; \textbf{chaque
pre/post/invariant relie au moins un vecteur} (une exigence invérifiable est une exigence
mal formulée). La provenance distingue les spécifications écrites avant le code — le cas
normal, exigé par R2 — des reverse-specs (§\ref{sec:reversespecs}).

\subsection{Révision et cycle de vie}

La révision s’incrémente à chaque changement du contrat. Une évolution qui rompt un
contrat existant ne modifie pas la spécification — elle en crée une nouvelle, l’ancienne
étant dépréciée avec renvoi : miroir exact du cycle de vie des blocs. Les appelants
dépendent du couple (identifiant, révision) ; la carte porte ce couple dans chaque entrée.

\subsection{Fixtures contractuelles partagées}

Quand plusieurs consommateurs se lient au même contrat externe — sortie d’un outil
système, format de fichier, API distante — les échantillons qui définissent ce contrat
appartiennent à la spécification, pas aux consommateurs. Le champ Fixtures référence un
jeu unique ; toute copie privée est un défaut de classe T5. Le mécanisme est né dans la
démonstration RoomEq : l’installateur scellé et l’application recodée partagent les
fixtures de parsing de pnputil — le défaut historique qui a survécu à trois versions
vivait précisément de la duplication.

\subsection{Taxonomie des défauts de contrat non écrit}\label{sec:taxonomie}

La classe de défauts que le protocole existe pour éliminer a un nom : \emph{contrat non
écrit}. La v0.4.0 en fixe la taxonomie (D19), née de l’analyse des défauts historiques de
la démonstration.

\begin{table}[htbp]
\centering\small
\caption{Taxonomie des défauts de contrat non écrit.}
\label{tab:taxonomie}
\begin{tabularx}{\textwidth}{@{}l Y Y@{}}
\toprule
\textbf{Classe} & \textbf{Définition} & \textbf{Verrou} \\
\midrule
T1 — valeur-seuil enfouie & Constante critique sans nom ni justification & constante
nommée + rationale dans la spec I \\
T2 — convention d’interface implicite & Structure ou ABI déclarée de mémoire & contrat
typé + invariant de taille testé \\
T3 — dépendance à l’hôte & Comportement dépendant de l’OS d’exécution & contrat
indépendant de l’hôte, union fixée \\
T4 — parsing de surface fragile & Analyse de libellés localisés ou volatils & fragments
indépendants des libellés + fixtures partagées \\
T5 — contrat dupliqué & Même connaissance codée en N exemplaires divergents & source
unique : spec I + fixtures \\
T6 — garantie non testée & Comportement critique sans vecteur de rejet & cas d’échec
nommé + test de rejet \\
\bottomrule
\end{tabularx}
\end{table}

Deux règles d’usage : tout défaut corrigé est classé dans cette taxonomie au journal des
décisions (la métrique « défauts ouverts par classe » du tableau de bord en dérive) ; la
taxonomie sert de liste de contrôle de rédaction — chaque bloc du gabarit verrouille au
moins une classe. Les occurrences de référence viennent de la démonstration : le seuil de
détection non justifié (T1+T6), la structure \code{SP\_DEVINFO\_DATA} de 424 octets au
lieu de 32 (T2, échec \code{0x6F8}), la sanitisation de noms dépendante de l’hôte (T3), le
parsing de pnputil sur libellés français (T4+T5, trois copies).

\section{Les trois couches : architecture et traçabilité}\label{sec:couches}

Le diff d’intention descend de I vers C, le verrou croisé relie C et B dans les deux sens,
et la couche B seule quitte le laboratoire. \textbf{Aucune couche n’est optionnelle} —
c’est une propriété normative : privé de I, le système perd son ancrage sémantique ; privé
de C, l’agent devient aveugle ; privé de B, il n’y a plus rien à exécuter. Toute
implémentation qui livrerait B sans I et C à jour viole R1 ; toute évolution qui
contournerait C pour éditer B directement viole R3 et la boucle.

\begin{table}[htbp]
\centering\small
\caption{Les trois couches.}
\label{tab:couches}
\begin{tabularx}{\textwidth}{@{}l l Y Y@{}}
\toprule
\textbf{Couche} & \textbf{Visibilité} & \textbf{Contenu} & \textbf{Rôle} \\
\midrule
I & Privée & Spécifications normalisées : contrats pré/post, invariants, cas d’échec,
vecteurs & Compression de l’intention \\
C & Privée, signée, chiffrée & Index des blocs, adresses, preuves, manifestes,
traçabilité & Unique interface de développement \\
B & Livrée & Code machine scellé, padding, checksums runtime & Exécution pure,
sémantiquement muette \\
\bottomrule
\end{tabularx}
\end{table}

La séparation croise une seconde dimension : la cible. La couche I est indépendante de la
cible — un contrat s’écrit une fois, ses types abstraits étant paramétrés par le profil.
La couche B est par profil. La couche C indexe l’ensemble : chaque entrée porte son
profil, et frag-bld organise les empreintes par profil. Quatre invariants structurels,
vérifiés à chaque cycle — leur violation est un refus de build : B n’est jamais nu (R1) ;
les adresses exposées sont stables (R7, par profil) ; la signature couvre la carte entière
(R9) ; B est auto-vérifiant au chargement (R10 renforcée).

\subsection{Traçabilité I $\leftrightarrow$ C $\leftrightarrow$ B $\leftrightarrow$
tests}\label{sec:traceabilite}

Chaque entrée de la carte référence quatre maillons : la spécification (identifiant +
révision), l’entrée d’index et le contrat embarqué, le(s) bloc(s) instancié(s) avec
empreinte(s) par profil, et les vecteurs exécutés avec le verdict du dernier cycle. Trois
propriétés sont contrôlées mécaniquement par la \textbf{commande d’audit de traçabilité} :
aucune spécification sans entrée de carte (l’exception visible étant la reverse-spec en
statut « prête ») ; aucun bloc scellé sans spécification référencée ; aucun vecteur
orphelin. Une chaîne incomplète est un refus de build — au même titre qu’un binaire nu.
La traçabilité est la condition de l’audit par des tiers : elle permet de répondre à
« quelle exigence ce bloc implémente-t-il, et quel test la défend ? » sans lecture du
binaire.

\section{La carte : fragmentation et budgets}\label{sec:carte}

\begin{sloppypar}
La carte est le cœur opérationnel du protocole : elle rend le binaire navigable, porte les
preuves et concentre la valeur intellectuelle. Elle est autodescriptive : un en-tête
normalisé consigne la version du protocole, la suite cryptographique~— hachage, signature,
chiffrement, chacun dans une liste blanche (D15)~— et les profils actifs. Sa spécification
détaillée fait l’objet de \code{FORMAT\_CARTE.md} — fragments, format d’entrée, verrou,
cycle de vie, export, re-cryptage —, mise à jour avec la présente révision : l’en-tête
porte \code{version=0.4.0} et chaque entrée de bloc porte le champ de traçabilité
\code{spec=<DOM>-<NN>:R}.
\end{sloppypar}

\paragraph{Norme de fragmentation.} La carte est fragmentée en quatre — frag-loc, frag-con,
frag-prf, frag-bld — et la v0.4.0 pose la norme : aucune opération de la boucle n’exige la
lecture de plus d’un fragment ; \textbf{aucune opération n’exige la lecture intégrale du
binaire} — le balayage est interdit (§\ref{sec:boucle}) ; l’audit complet de la carte est
une opération exceptionnelle, non budgétée, journalisée.

\paragraph{Budgets par type d’opération.}

\begin{table}[htbp]
\centering\small
\caption{Budgets en tokens par opération (localiser amendé à 64 tokens par D31 après la mesure M1 — le gabarit canonique v0.4.0 mesurait 43 tokens en cl100k\_base ; autres budgets inchangés depuis la v0.3).}
\label{tab:budgets}
\begin{tabularx}{\textwidth}{@{}Y Y l@{}}
\toprule
\textbf{Opération} & \textbf{Données lues} & \textbf{Budget} \\
\midrule
Localiser un bloc & frag-loc — une entrée & 64 tokens max (D31) \\
Éditer simplement un bloc & frag-loc + frag-con + notation & 3 000 tokens max \\
Ajouter un bloc & frag-loc + frag-con + seuils frag-prf & 6 000 tokens max \\
Instancier un profil nouveau & I partagée + frag-con + notation du profil & 6 000 tokens
max par bloc \\
Audit complet (rare) & carte entière déchiffrée & non budgété \\
\bottomrule
\end{tabularx}
\end{table}

Indicateurs de santé : précision de localisation (cible 95 \% dès M0), coût moyen d’une
édition simple (cible 5 000 tokens), taux de régression post-patch (cible zéro).
Instrumentés depuis M1 : la mesure effective (51 / 719 / 1 088 tokens, D31) a corrigé
le budget de localisation à 64 et valide les deux autres ; le tableau de bord les
affiche comme métriques mesurées, avec leur écart.
   % 1 Vision · 2 Couche I + taxonomie · 3 Couches + traçabilité · 4 Carte
% =====================================================================
%  CH 5-7 : Règles · Boucle de mise à jour · Vérification
% =====================================================================
\section{Les règles du protocole}\label{sec:regles}

R1 à R7 proviennent de la v0.1 ; R8 à R10 de la v0.2 ; \textbf{R10 est renforcée et R11
ajoutée par la v0.4.0}. Les règles sont peu nombreuses, vérifiables mécaniquement et
sanctionnées par des effets concrets — refus de build, refus de commit ou blocage de
session — plutôt que par des recommandations. Aucune règle existante n’est affaiblie : les
énoncés de R1 à R9 sont repris à l’identique, celui de R10 est un sur-ensemble strict de
la version v0.3.

\begin{table}[htbp]
\centering\small
\caption{Les onze règles du protocole v0.4.0.}
\label{tab:regles}
\begin{tabularx}{\textwidth}{@{}l l Y l@{}}
\toprule
\textbf{\#} & \textbf{Règle} & \textbf{Énoncé} & \textbf{Sanction} \\
\midrule
R1 & Jamais de binaire nu & Tout artefact B est accompagné de I et C à jour & build refusé \\
R2 & Intention d’abord & Toute modification commence par un diff sur I & édition rejetée \\
R3 & Balisage obligatoire & Chaque bloc est décrit dans C & build refusé \\
R4 & Patch à contrat stable & Le remplacement préserve signatures et contrats ; les
appelants ne sont jamais invalidés & échec de vérification \\
R5 & Vérifie-puis-écrit & Aucun commit sans oracle, vecteurs dorés et round-trip &
commit impossible \\
R6 & Cristallisation & Dix vérifications réussies rendent le bloc immuable ; le modifier
exige la procédure lourde & nouvelle version du bloc \\
R7 & Adresses stables & Les points d’entrée sont garantis par padding ; patcher ne
déplace rien & re-liage forcé — interdit \\
R8 & Verrou croisé & Empreintes des blocs et signature de la carte vérifiés avant chaque
édition & session bloquée \\
R9 & Carte signée obligatoire & L’agent refuse toute carte non signée & refus de
démarrage \\
\textbf{R10} & Échec contrôlé intégré \emph{(renforcée)} & Au lancement — avant toute
entrée dans le chemin d’exécution — l’artefact vérifie (a) les checksums de ses blocs
scellés, (b) la signature du manifeste de build couvrant code et données embarquées, (c)
la monotonie de version (anti-rollback : tout état antérieur à la version minimale
acceptée est refusé), (d) la conformité minimale de l’environnement au profil. Tout écart
= arrêt propre immédiat ; aucun chemin d’exécution dégradé n’existe \emph{(fail-closed
explicite)} & arrêt avant tout effet \\
\textbf{R11} & Optimisation mesurée \emph{(nouvelle)} & Toute optimisation exige un hot
path identifié par mesure, une mesure avant/après consignée au journal d’optimisation et
un contrat stable (R4). L’optimisation prématurée — sans hot path identifié ni mesure —
est interdite : la demande est refusée et consignée & refus de commit \\
\bottomrule
\end{tabularx}
\end{table}

Quatre précisions. R4 : un bloc remplaçant honore les contrats publiés de son prédécesseur
— mêmes signatures, mêmes pré et post-conditions, mêmes invariants visibles — condition de
possibilité du patch à chaud. R7 s’entend \textbf{par profil} : les adresses sont stables
au sein du binaire de chaque profil. R1, précisée par D16, interdit de livrer la carte
avec l’artefact — pas de détenir et transférer son propre état : c’est l’objet de l’export
de continuité. R10 renforcée (D20) étend le fail-closed en un contrôle intégré à quatre
volets ; la vérification est strictement \textbf{locale} (aucune dépendance réseau au
chargement) et la politique anti-rollback précise (registre des versions acceptées,
stockage du plancher) appartient au profil de cible.

\section{La boucle de mise à jour}\label{sec:boucle}

La boucle est la seule voie par laquelle le logiciel évolue : sept étapes en deux phases —
l’édition, puis la validation et le commit — séparées par un point de non-retour : rien
n’est écrit dans le binaire avant que l’oracle n’ait rendu son verdict.

\begin{table}[htbp]
\centering\small
\caption{La boucle de mise à jour en sept étapes.}
\label{tab:boucle}
\begin{tabularx}{\textwidth}{@{}l l Y l@{}}
\toprule
\textbf{N} & \textbf{Étape} & \textbf{Entrées $\to$ sorties} & \textbf{Règles} \\
\midrule
1 & Diff d’intention & requête de changement, couche I $\to$ diff minimal sur I & R2 \\
2 & \textbf{Localisation par intention} & frag-loc, diff I $\to$ blocs impactés & R3 \\
3 & Verrou croisé & frag-bld, binaire $\to$ verdict go / no-go & R8, R9 \\
4 & Régénération ciblée & frag-con, notation $\to$ octets candidats & R4, R7, R11 \\
5 & Vérification & candidats, oracle, vecteurs $\to$ verdict, journal & R5 \\
6 & Commit atomique & verdict favorable $\to$ patch B, carte, scellage & R1 \\
7 & Cristallisation & compteur de vérifications $\to$ bloc immuable & R6 \\
\bottomrule
\end{tabularx}
\end{table}

\paragraph{Localisation par intention (norme v0.4.0).} L’étape 2 porte désormais son nom
entier : la localisation est \textbf{dérivée du diff d’intention} par l’index de la carte
— jamais par balayage du binaire, jamais par recherche de motif dans les octets. Un diff
qui ne localise aucun bloc est une erreur de spécification à corriger dans I, pas une
licence de scan ; un balayage serait une violation du budget (§\ref{sec:carte}) et de R3.
La précision de localisation — bon bloc identifié du premier coup — reste l’indicateur de
santé mesuré dès M0 (cible 95 \%).

\paragraph{Patch à chaud du profil primaire.} Le mécanisme exact est spécifié par chaque
profil ; pour x86-64 : le nouveau bloc est d’abord écrit dans une zone de croissance
libre, puis un trampoline de 12 à 16 octets est installé à l’ancien point d’entrée et
bascule par une écriture alignée unique, atomique pour le processeur. L’ancien trampoline
est conservé : le rollback copy-on-write consiste à rebasculer vers lui. Les adresses
stables garantissent qu’aucun appelant ne référence l’intérieur d’un bloc : seuls les
points d’entrée~— promis par R7~— sont trampolinés, et leur table est close et vérifiée
par le verrou croisé. Le critère M2 : remplacer le bloc d’écho d’un serveur TCP pendant
qu’il traite des connexions, sans interruption, avec rollback démontré~— \textbf{démontré
sur machine réelle le 27 octobre 2026} (D43) : la boucle complète (bascule pendant le
trafic, vérification memcmp, retour arrière, re-vérification) s’est terminée par le
constat « DÉMO OK » du journal d’exécution, au terme de la campagne relatée en
§\ref{sec:terrain}.

\section{La stratégie de vérification}\label{sec:verification}

Supprimer le compilateur, c’est perdre les filets de sécurité qu’il tendait gratuitement.
La stratégie de vérification reconstruit une confiance équivalente sans réintroduire de
couche de traduction dans l’artefact.

\paragraph{L’oracle éphémère.} Implémentation de référence de chaque bloc — triviale,
volontairement lente, dérivée de la couche I, générée au début de chaque session puis
détruite. Le test est différentiel par propriétés : l’oracle et le bloc machine reçoivent
des entrées aléatoires contraintes par les pré-conditions, leurs sorties sont comparées
octet par octet ; volume de référence : douze mille cas par bloc critique — les campagnes
effectives sont consignées au journal (76 010 cas en M1, 50 533 en M2, 86 pour la
démonstration graphique). L’oracle dérive
de I et non de B ; il est régénéré à chaque session ; sa régénération coûte moins que le
diagnostic d’une corruption post-scellage. Dérivé de la couche I — indépendante de la
cible — le \textbf{même} oracle valide les blocs de tous les profils : la portabilité ne
divise pas le coût de vérification, elle le mutualise.

\paragraph{Le round-trip de notation.} Chaque bloc assemblé est désassemblé et comparé à
la notation d’origine ; tout désaccord — instruction, opérande, taille d’encodage — est un
échec de build. La notation n’est pas un compilateur : aucune optimisation, aucune
réécriture ne s’intercale entre l’écriture et les octets (principe WYSIWYG).

\paragraph{La cristallisation.} Plus un bloc fait ses preuves, moins il coûte à vérifier.
Chaîne de confiance du scellage : verdict de l’oracle $\to$ vecteurs dorés hachés dans
frag-prf $\to$ empreinte du bloc dans frag-bld $\to$ signature de la carte entière $\to$
immuabilité à dix vérifications.

\subsection{Vecteurs de vérification}\label{sec:vecteurs}

Trois espèces : \textbf{vecteurs dorés} (entrées/sorties de référence, conservés hachés) ;
\textbf{propriétés différentielles} (invariants testés par l’oracle sur entrées
contraintes) ; \textbf{tests de rejet} (chaque cas d’échec du contrat, matérialisé par un
test qui exige la sanction nommée). Exigence de couverture : chaque exigence numérotée du
contrat et chaque cas d’échec relie au moins un vecteur — la commande d’audit de
traçabilité vérifie ce maillon.


\subsection{La validation terrain — la boucle d'évidence}\label{sec:terrain}

La v0.5.0 élève la validation terrain au rang de couche de vérification normée (D43). Le
principe est né d'un constat structurel : l'oracle différentiel exécute le binaire contre un
\textbf{modèle} de l'environnement — chargeur, ABI, sémantique des API, protections mémoire.
Un binaire peut donc être exact au sens du round-trip, intègre au sens de la porte R10,
conforme au sens de la suite du profil, vert au sens de l'oracle — et mourir au premier appel
sur machine réelle, parce que le modèle était plus pauvre que la réalité. Les jalons M1 et
M2 en ont fourni la preuve expérimentale : huit classes de défauts ainsi invisibles
(§\ref{sec:taxonomieX}) n'ont été révélées que par l'exécution terrain.

La boucle d'évidence comporte cinq éléments normés :

\begin{enumerate}
\item \textbf{Le journal d'exécution} — l'artefact livré écrit ses preuves au fil de l'eau
(fichier à côté de l'exécutable, ligne par ligne). La règle d'honnêteté est absolue : une
ligne n'est écrite qu'après la vérification de l'étape qu'elle atteste (memcmp exact) — une
ligne présente vaut preuve ; la dernière ligne visible localise toute mortalité à
l'intervalle près.
\item \textbf{La lecture asymétrique} — le terrain retourne son verdict par un canal à
faible bande passante : le journal (éventuellement complété d'une capture). L'atelier
n'exige du terrain que ce document, jamais un diagnostic : la charge d'analyse reste à
l'agent.
\item \textbf{L'instrumentation contractualisée} — quand l'intervalle de mortalité reste
aveugle, l'itération suivante embarque des sondes de trace aux frontières (D40) : elles
journalisent des faits (positions d'exécution, retours vérifiés), jamais un succès non
prouvé ; leur retrait après constat vert est promis dans la spécification (D43 a honoré la
promesse) et vérifié par l'oracle — le journal nominal redevient exactement sa séquence de
preuves, toute ligne parasite est un échec.
\item \textbf{Le diff chirurgical} — chaque itération terrain ne modifie que les blocs
révisés ; l'audit d'itération vérifie mécaniquement qu'aucun octet hors des zones
autorisées (blocs révisés, chaînes, version, sommes) n'a bougé — l'héritage scellé reste
octet-identique (R7).
\item \textbf{L'archive brute} — chaque journal terrain est archivé tel quel (octets
d'origine, encodage et fins de ligne compris) avec son empreinte : la preuve est le
document, pas son résumé.
\end{enumerate}

Le protocole est ensuite économique : huit lectures de journal ont suffi à nommer huit
classes de défauts en deux jours (demo.1 à demo.8), et le dernier diagnostic a été obtenu
sans run supplémentaire — par lecture statique des octets appuyée sur les preuves
accumulées (D42). La contrepartie est la garantie de progression : chaque défaut nommé
devient un vecteur de falsification permanent de l'oracle — le modèle d'environnement
s'enrichit de chaque lecture terrain et ne peut plus régresser en silence.

\subsection{Taxonomie étendue — les défauts invisibles à l'émulation (X1--X8)}\label{sec:taxonomieX}

Les classes T1--T6 (§\ref{sec:taxonomie}) classent les défauts de contrat non écrit :
lisibles dans la spécification, corrigeables en écrivant le contrat. La campagne M2 a
révélé une seconde famille, structurellement différente : des défauts dont ni la
spécification, ni la notation, ni le round-trip, ni l'oracle ne parlent — parce que le
modèle d'environnement de l'oracle était, sur un point précis, plus pauvre que le système
réel. La v0.5.0 en fixe la taxonomie (D35--D43) :

\begin{table}[htbp]
\centering\footnotesize
\caption{Taxonomie des défauts invisibles à l'émulation (classes X).}
\label{tab:taxonomieX}
\begin{tabularx}{\textwidth}{@{}l l Y Y Y@{}}
\toprule
\textbf{Classe} & \textbf{Déc.} & \textbf{Modèle défaillant} & \textbf{Défaut de
référence} & \textbf{Garde permanente} \\
\midrule
X1 & D35 & sémantique du chargeur & \code{FirstThunk} $\neq$ slots appelés : le chargeur
résout ailleurs que le code n'appelle & verrou « chargeur réel » : FT = slots appelés \\
X2 & D36 & lecture de structure & \code{RECT.top} lu comme largeur : spec juste, octets
faux — bloc jamais exécuté, shim muet & coquille exécutée + invariant géométrique \\
X3 & D37 & canal de vérité & jumeau silencieux : le fichier vivait, l'écran restait vide —
défaut dans la spécification & miroir universel, les deux canaux contrôlés \\
X4 & D38 & ABI caller-saved & volatils relus après l'appel : un retour de syscall écrit
RFLAGS dans r11 & clobber des volatiles à chaque retour simulé \\
X5 & D39 & protections mémoire & phase de patch \code{PAGE\_READWRITE} sans execute : DEP
permanent x64 & le shim applique les permissions réelles \\
X6 & D40 & écart émulation/terrain & défaut latent qu'aucune exécution n'avait déclenché :
audit exhaustif vert & sondes aux frontières, instrumentation contractualisée \\
X7 & D41 & obligations ABI de l'appelant & zone d'accueil msx64 absente : l'appelé écrit
sur l'adresse de retour & shim \code{/homeparams} + audit systémique des sites \\
X8 & D42 & état hérité & section critique exigée par un bloc scellé, initialisée par un
bloc orchestrateur retiré & shims à état + exécution réelle du consommateur \\
\bottomrule
\end{tabularx}
\end{table}

Trois règles d'usage, miroir de celles de la taxonomie T. Premièrement, tout défaut terrain
corrigé est classé T ou X au journal des décisions — la métrique « défauts ouverts par
classe » du tableau de bord couvre désormais les deux familles. Deuxièmement, chaque classe
X est gardée par une falsification : réintroduire le défaut doit faire échouer l'oracle —
le modèle s'enrichit mais ne régresse jamais en silence. Troisièmement, la taxonomie est
ouverte : une neuvième classe sera nommée X9 avec sa décision, jamais dissimulée dans un
correctif muet.

La lecture académique : un oracle différentiel valide la conformité d'un artefact à un
modèle ; la validation terrain valide le modèle lui-même. Les deux couches sont
complémentaires, non redondantes — l'oracle couvre l'espace des entrées à coût nul, le
terrain couvre l'espace des environnements au coût d'un run. Le protocole exige les deux :
R5 (vérifie-puis-écrit) pour la première, la boucle d'évidence pour la seconde.
\subsection{Suite de conformité par profil}\label{sec:conformite}

Chaque profil embarque une \textbf{suite de conformité} vérifiant ses sept éléments
constitutifs (D22) :

\begin{table}[htbp]
\centering\small
\caption{Suite de conformité d’un profil : les sept contrôles.}
\label{tab:suite}
\begin{tabularx}{\textwidth}{@{}l l Y@{}}
\toprule
\textbf{\#} & \textbf{Élément du profil} & \textbf{Contrôle de la suite} \\
\midrule
1 & ISA et types & tailles, alignements, boutisme constatés = déclarés \\
2 & ABI d’appel & appels et retours conformes ; volatilité ; shadow space ou zone rouge \\
3 & Format d’exécutable & en-têtes valides ; protections constatées (DEP/NX, ASLR/PIE,
CFI/CFG) \\
4 & Interface système & résolution des imports ou syscalls ; contrat d’erreurs \\
5 & Patch à chaud & bascule atomique et invalidation du cache d’instructions démontrées \\
6 & Notation & round-trip 100 \% sur le corpus d’instructions du profil \\
7 & Granularité et padding & adresses stables après patch (R7) \\
\bottomrule
\end{tabularx}
\end{table}

Un profil sans suite verte n’est pas déclarable supporté — un profil incomplet est refusé
exactement comme un build sans carte (D14). La suite est un livrable obligatoire de
l’instanciation, exécutée à chaque cristallisation de version rc et finale, et le critère
d’acceptation des jalons M6 et M7.

\subsection{Parité inter-profils à contrat constant}\label{sec:parite}

Deux profils sont \textbf{paritaires} pour le contrat $\sigma$ si et seulement si :
(1) $\sigma$ est instancié sur les deux depuis la même spécification I, à la
paramétrisation près ; (2) l’oracle — dérivé de I, mutualisé — est vert sur les deux ;
(3) les verdicts des vecteurs de $\sigma$ sont équivalents ; (4) les différences de profil
sont documentées dans frag-con de chaque bloc. La parité se \textbf{vérifie par
exécution}, jamais par déduction : c’est le critère de sortie de M6 (inter-OS) et M7
(inter-ISA). Aucune couche d’abstraction n’est introduite pour l’obtenir — un bytecode
portable trahirait la fonction objectif Performance (rejet consigné, §\ref{sec:rejets}).
   % 5 Règles · 6 Boucle · 7 Vérification
% =====================================================================
%  CH 8-11 : Pré-versions · Reverse-specs · Optimisation · Profils
% =====================================================================
\section{Cristallisation et pré-versions}\label{sec:preversions}

La v0.3 connaissait les balises de version sans en fixer la sémantique ; la démonstration
RoomEq a inventé en cours de route la pré-version \code{-dev.N} pour distinguer « vérifié
par l’oracle » de « validé sur le terrain ». La v0.4.0 formalise le schéma complet (D23) :

\begin{table}[htbp]
\centering\small
\caption{Schéma de pré-versions de cristallisation.}
\label{tab:preversions}
\begin{tabularx}{\textwidth}{@{}l l Y@{}}
\toprule
\textbf{Segment} & \textbf{Signification} & \textbf{Conditions d’ouverture} \\
\midrule
\code{X.Y.Z-dev.N} & pré-version de développement & oracle vert (R5) ; construction
0 erreur / 0 avertissement ; jamais validée terrain ; usage interne \\
\code{X.Y.Z-rc.N} & candidate à la validation & dev + suite de conformité du profil
verte (§\ref{sec:conformite}) + critères du jalon courant ; soumise à la validation
terrain \\
\code{X.Y.Z} & version finale & validation terrain consignée au journal des décisions ;
balise immuable \\
\bottomrule
\end{tabularx}
\end{table}

Trois règles gouvernent le schéma. \textbf{Monotonie} : pour un même triplet,
\code{dev.N} $<$ \code{rc.N} $<$ version finale ; la version minimale acceptée alimente
le contrôle anti-rollback de R10. \textbf{Immuabilité} : une balise posée n’est jamais
déplacée ni supprimée ; un défaut découvert après cristallisation se corrige en
\code{X.Y.Z+1} à contrat stable (R4) — l’échec est honoré, pas effacé.
\textbf{Progression de contrat} : une évolution qui rompt un contrat publié passe par
\code{X.(Y+1).0} avec révision explicite de la spécification concernée. Le schéma
généralise la séquence vécue par RoomEq (\code{v3.0.0-dev.1} $\to$ \code{v3.0.0} $\to$
\code{v3.0.1} $\to$ \code{v3.0.2} : deux échecs d’intégration continue honorés puis
corrigés en patchs à contrat stable), qui demeure l’occurrence de référence. Le schéma
a depuis été vécu deux fois par le cœur natif : les balises finales \code{0.2.0} (M1,
D32) et \code{0.3.0} (M2, D43) sont posées, chacune après validation terrain consignée
— la séquence rc.1 $\to$ rc.4 de M2 (quatre générations, chacune sur défaut terrain
nommé et corrigé) demeure l’occurrence de référence de la boucle d’évidence
(§\ref{sec:terrain}).

\section{Binaires scellés hérités : reverse-specs}\label{sec:reversespecs}

Un projet agent-first hérite presque toujours d’un patrimoine : binaires scellés, validés
sur le terrain, dont les contrats ne sont écrits nulle part. La v0.3 ne disait rien de ce
cas ; la démonstration y a répondu par invention (spécifications de l’installateur scellé
v2.2.2, huit modules extraits par rétro-ingénierie). La v0.4.0 officialise la pratique
(D23).

\begin{sloppypar}
\paragraph{Définition.} Une \emph{reverse-spec} est une spécification I extraite d’un
artefact scellé existant : le gabarit s’applique intégralement, avec deux particularités.
La \textbf{provenance} est \code{reverse-spec de <binaire scellé @ tag>} — elle identifie
l’artefact source et sa balise. Le \textbf{statut} est « prête » tant que la spécification
n’a pas été confrontée au comportement réel (fixtures issues de l’artefact, tests de
discrimination) ; il passe à « validée » lorsque les vecteurs ont tranché. Une
reverse-spec ne modifie jamais l’artefact scellé : elle le documente.
\end{sloppypar}

\paragraph{Règles.} Premièrement, le recodage d’un binaire scellé ne peut partir que de sa
reverse-spec \emph{validée} — jamais du code relu en séance, jamais d’une mémoire de
session. Deuxièmement, une reverse-spec « prête » figure dans la traçabilité avec son
statut visible : l’exception est nommée, pas silencieuse. Troisièmement, la
rétro-ingénierie se fait par lecture ligne à ligne du code source scellé — méthode
déclarée dans la spécification, reproductible par un tiers. La reverse-spec étend la
couche I au patrimoine sans l’exiger de recoder : figer la connaissance avant de décider.

\section{Optimisation mesurée et journal d’optimisation}\label{sec:optimisation}

La fonction objectif place l’optimisation au troisième rang et l’assortit d’une condition :
descente à l’ISA sur hot paths identifiés par profilage. La v0.4.0 en tire l’interdiction
normative (R11) et son instrument, le \textbf{journal d’optimisation} (D21) :
l’optimisation sans hot path identifié ni mesure est prématurée, donc interdite ; sa
demande est refusée et consignée — le refus motivé fait partie de la gouvernance, comme
les alternatives rejetées du journal des décisions.

\begin{lstlisting}[caption={Format normalisé d’une entrée du journal
d’optimisation.}, label={lst:opt}]
[OPT-NN] date . bloc/chemin . decideur
  hot path       : <identification + preuve de mesure>
  mesure avant   : <valeur, conditions, materiel>
  changement     : <description> . contrat stable : oui/non (R4)
  mesure apres   : <valeur, conditions identiques>
  gain           : <x ou %> . decision : appliquee | refusee (motif)
  decision liee  : Dxx
\end{lstlisting}

Le journal est versionné avec la couche I ; toute optimisation appliquée doit précéder le
commit du bloc optimisé (R11 sanctionne le commit, pas seulement l’intention). L’entrée de
référence est l’occurrence RoomEq : la demande d’« optimisation maximale » a été honorée
par la mesure — chaîne d’égalisation à 152$\times$ le temps réel, zéro allocation par
bloc — puis refusée comme transformation, faute de problème de performance à corriger.

\section{Profils de cible}\label{sec:profils}

Le cœur du protocole — couches, règles, boucle, vérification — est indépendant de la cible
par construction ; tout ce qui dépend de la cible est regroupé dans un \textbf{profil}.
L’universalité reste structurelle : instancier une plateforme nouvelle revient à écrire un
profil — ses sept éléments \emph{et sa suite de conformité} (§\ref{sec:conformite}) —
jamais à modifier le protocole.

Un profil définit sept éléments : (1) l’ISA et ses types ; (2) l’ABI d’appel ; (3) le
format d’exécutable et les protections à honorer ; (4) l’interface système autorisée ;
(5) le mécanisme de patch à chaud natif ; (6) le moteur de notation et son vérificateur
de round-trip ; (7) la granularité des blocs et le padding des adresses stables. Un
huitième livrable s’ajoute en v0.4.0 : la \textbf{politique anti-rollback et de
vérification au chargement} propre au profil (déclinaison de R10).

\begin{table}[htbp]
\centering\small
\caption{Profils de cible déclarés.}
\label{tab:profils}
\begin{tabularx}{\textwidth}{@{}l l l Y@{}}
\toprule
\textbf{Profil} & \textbf{Format} & \textbf{Interface} & \textbf{Statut} \\
\midrule
\code{x86-64-win11-pe} & PE32+ & IAT (kernel32, ws2\_32, Schannel) & primaire — jalons
M0–M5 \\
\code{x86-64-linux-elf64} & ELF64 & syscalls + libc & M6 — inter-OS \\
\code{riscv64-linux-elf64} & ELF64 & syscalls + libc & M7 — inter-ISA \\
\code{aarch64-macos-macho} & Mach-O & libSystem & réservé \\
\code{wasm32-wasi} & Wasm & WASI & réservé, étude \\
\bottomrule
\end{tabularx}
\end{table}

\paragraph{Profil primaire, ABI Microsoft x64.} Arguments en RCX, RDX, R8, R9 puis la
pile ; retour en RAX ; shadow space de 32 octets par appelant ; volatils RAX, RCX, RDX,
R8–R11, XMM0–XMM5 ; RSP aligné 16 ; \textbf{aucune zone rouge} — contrairement à System
V, stocker sous RSP sans allocation préalable est interdit. R10 et R11 sont les registres
de travail internes ; aucune donnée vivante ne séjourne entre RSP et RSP+32 au moment
d’un appel ; tout bloc appelant alloue son cadre avant d’y toucher.

\paragraph{Différences entre profils.} Le passage Windows $\to$ Linux \textbf{inverse} le
piège de la zone rouge (interdite sur msx64, permise sur SysV pour une fonction feuille) ;
le passage x86-64 $\to$ RISC-V change la nature du patch à chaud : sur une ISA à
instructions de taille fixe, la bascule passe par un registre scratch chargé de l’adresse
cible, suivi de \code{fence.i} — le profil spécifie la séquence exacte et son atomicité.
   % 8 Pré-versions · 9 Reverse-specs · 10 Optimisation · 11 Profils
% =====================================================================
%  CH 12-16 : Gouvernance · Menaces et limites · Jalons · Correspondance · Décisions
% =====================================================================
\section{Gouvernance}\label{sec:gouvernance}

\subsection{Journal des décisions — format normalisé}\label{sec:journal}

Le protocole ne se révisant jamais en silence, chaque décision structurante est consignée
dans un journal. La v0.4.0 normalise le format (D24) :

\begin{table}[htbp]
\centering\small
\caption{Format normalisé d'une entrée du journal des décisions.}
\label{tab:formatdecisions}
\begin{tabularx}{\textwidth}{@{}l Y@{}}
\toprule
\textbf{Champ} & \textbf{Exigence} \\
\midrule
ID & \code{DNNN}, séquentiel, jamais réutilisé \\
Date & date de consignation \\
Statut & active · amendée par Dxx · invalidée par Dxx \\
Contexte & exigence, incident ou évaluation qui motive la décision \\
Choix retenu & décision en une phrase, rattachée à la fonction objectif \\
Alternatives rejetées & au moins une, avec motif \\
Amendements & référence explicite à la décision remplacée, le cas échéant \\
Développement & argumentaire détaillé si la décision engage l'architecture \\
\bottomrule
\end{tabularx}
\end{table}

Trois règles : le journal est \textbf{append-only} (une entrée ne s'efface jamais) ; une
révision \textbf{invalide explicitement} ce qu'elle remplace ; aucune entrée n'est
rétrodatée. La conformité du journal — présence des champs, unicité des ID, chaînage des
amendements — est contrôlée mécaniquement par l'audit de gouvernance.

\subsection{Tableau de bord de conformité}\label{sec:tableaudebord}

La v0.4.0 unifie les indicateurs de santé en un \textbf{tableau de bord} (D24), généré par
la commande d'audit à chaque commit et à chaque cristallisation, archivé avec la carte :

\begin{table}[htbp]
\centering\small
\caption{Tableau de bord de conformité : métriques, cibles, sources.}
\label{tab:dashboard}
\begin{tabularx}{\textwidth}{@{}Y l Y@{}}
\toprule
\textbf{Métrique} & \textbf{Cible} & \textbf{Source de mesure} \\
\midrule
Blocs avec spécification I à jour & 100 \% & audit de traçabilité \\
Chaînes de traçabilité complètes I $\leftrightarrow$ C $\leftrightarrow$ B
$\leftrightarrow$ tests & 100 \% & audit de traçabilité \\
Construction & 0 / 0 & build \\
Tests différentiels & 100 \% verts & oracle \\
Empreintes des données embarquées & conformes & hachage vs carte \\
Défauts ouverts par classe T1–T6 & 0 & journal des décisions \\
Optimisations sans mesure (R11) & 0 & journal d'optimisation \\
Budgets tokens (localiser / éditer / ajouter) & $\leq$ 64 / 3 000 / 6 000 &
instrumentation — mesurés M1 : 51 / 719 / 1 088 (D31) \\
Suite de conformité par profil actif & verte & suite du profil \\
Changements hors journal des décisions & 0 & diff vs journal \\
Pré-versions correctement balisées (dev / rc / final) & 100 \% & registre des balises \\
\bottomrule
\end{tabularx}
\end{table}

Le tableau de bord est la contrepartie mesurable de la gouvernance : chaque affirmation du
protocole s'y retrouve sous forme de métrique avec sa cible et sa source — aucune ne
s'appuie sur une déclaration non vérifiable. Les métriques dépendantes de l'outillage
affichent leur statut « cible » jusqu'à leur premier mesurage.

\subsection{Continuité et anti-enfermement}

L'export de continuité (D16) demeure l'unique forme de transfert de l'état de
développement : conteneur chiffré et signé regroupant I, B, C et un manifeste
autodescriptif ; oracle régénéré depuis I à l'arrivée, jamais exporté ; import en cinq
étapes dont l'échec d'une seule est terminal. Sauvegarde 3-2-1 chiffrée et escrow de la
clé maîtresse inchangés.

\paragraph{Le coffre-fort des phrases maîtresses (D75).} Le composant qui centralise
toutes les phrases maîtresses d'un projet — et de tout projet ultérieur conduit sous le
protocole — dans un conteneur unique chiffré TRIADE-CF1 (\code{coffre\_fort.py}) :
phrase de contrôle détenue par le seul développeur (jamais stockée), Argon2id aux
planchers de la carte au repos puis HKDF à contexte dédié, XChaCha20-Poly1305 avec
l'en-tête clair pour AAD, écriture atomique 0600, autotest de falsification exigé à
chaque installation. Le conteneur chiffré s'exporte tel quel vers le coffre Drive du
projet et le coffre personnel du développeur — son vol relève de la menace « sans la
phrase, indéchiffrable ». Après une réinitialisation du poste, l'import du conteneur
puis \code{obtenir --vers-fichier} restaurent la phrase opérationnelle sans cérémonie
de régénération. Discipline inchangée : chaque cérémonie de clés (D27/D30) alimente le
coffre dans la même session, la phrase de contrôle est affichée une seule fois, et la
graine Ed25519 reste hors coffre (D30) — format normatif : FORMAT\_CARTE § 11.

\section{Menaces et limites}\label{sec:menaces}

Le principe directeur demeure : séparer l'obfuscation, qui protège la propriété
intellectuelle, du durcissement, qui protège l'exécution. Conformément au principe de
Kerckhoffs, la sécurité ne repose jamais sur le secret de la carte : elle survivrait à sa
divulgation.

\subsection{Modèle de menace (étendu en v0.4.0, D20)}

\begin{table}[htbp]
\centering\small
\caption{Modèle de menace — ajouts v0.4.0 en gras ; ajout v0.5.1 en gras (D75).}
\label{tab:menaces}
\begin{tabularx}{\textwidth}{@{}Y Y Y@{}}
\toprule
\textbf{Menace} & \textbf{Contre-mesure} & \textbf{Résiduel} \\
\midrule
Rétro-ingénierie de l'IP & carte privée ; binaire sémantiquement muet ; variantes par
build & temps acheté — s'érode avec les LLM de décompilation \\
Empoisonnement de carte & signature ; refus de carte non signée ; verrou croisé (R8, R9)
& détecté — session bloquée \\
Corruption de carte (accidentelle) & AEAD par fragment + signature + verrou & détecté au
chargement \\
Fuite de la carte & chiffrement ; sharding ; escrow & perte limitée à un fragment \\
Vol de l'export de continuité & chiffrement intégral + signature & indéchiffrable sans
la phrase \\
\textbf{Vol du coffre-fort des phrases} & chiffrement intégral (phrase de contrôle +
Argon2id, contexte HKDF dédié) ; AEAD ; écriture atomique & indéchiffrable sans la
phrase (D75) \\
\textbf{Binaire altéré} & checksums par bloc + signature du manifeste au lancement
(R10 a, b) & détecté — arrêt avant tout effet \\
\textbf{Rollback non autorisé} & monotonie de version signée, plancher anti-rollback
(R10 c) & détecté — refus \\
\textbf{Agent défaillant} & R2 ; R5 ; oracle dérivé de I ; budgets & l'erreur ne devient
jamais permanente sans vérification \\
\textbf{Données embarquées modifiées} & empreintes dans la carte ; vérification au
lancement (R10 b) & détecté — arrêt propre \\
\textbf{Environnement non conforme} & contrôles de conformité au profil (R10 d) ;
fail-closed & refus propre avant exécution \\
Fuzzing boîte noire & contrats, oracle, tests par propriétés & indépendant du secret de
la carte \\
Redistribution frauduleuse & empreintes par build — canary trap & fuite traçable \\
Instrumentation dynamique & anti-debug léger par mesure de temps & assumé \\
\bottomrule
\end{tabularx}
\end{table}

Les six menaces en gras sont les additions de la v0.4.0, le vol du coffre-fort (D75)
celle de la v0.5.1 ; toutes sont neutralisées par des mécanismes \textbf{locaux et
vérifiables} — la contre-mesure anti-rollback ne dépend d'aucun service en ligne (le
chargement doit échouer proprement hors ligne), et les vérifications de R10 s'exécutent
avant toute entrée dans le chemin d'exécution.

\subsection{Limites assumées}\label{sec:limites}

La rétro-ingénierie est un objectif d'économie, pas d'impossibilité : le coût de l'analyse
d'un binaire strippé, sans structure, dont chaque exemplaire diffère, reste supérieur à sa
valeur — mais les modèles de langage appliqués à la décompilation progressent, et
l'avantage doit être réévalué à chaque génération. L'anti-debug reste léger et assumé.
La furtivité face aux EDR est hors périmètre : un outil légitime par construction, dont
la transparence vis-à-vis du système est totale. Enfin, la gouvernance du protocole a été
validée sur un cas unique (N = 1) — l'étude de cas RoomEq, dont les limites sont
documentées dans l'édition académique (D18) : la généralité reste à établir par
réplication, le cœur natif par les jalons M0 à M7.

\section{Critères d'acceptation des jalons M0 à M7}\label{sec:jalons}

Huit jalons, chacun avec des critères de sortie mesurables (détail dans
\code{PLAN\_MILESTONES.md}). \textbf{Conditions transverses} (v0.4.0) : aucun jalon n'est
franchi sans oracle vert, suite de conformité verte pour les éléments couverts, tableau de
bord conforme et carte signée.

\begin{table}[htbp]
\centering\small
\caption{Critères d'acceptation des jalons.}
\label{tab:jalons}
\begin{tabularx}{\textwidth}{@{}l Y >{\raggedright\arraybackslash}p{2.6cm}@{}}
\toprule
\textbf{Jalon} & \textbf{Critères de sortie mesurables} & \textbf{Statut} \\
\midrule
M0 — chaîne artefact & exécution sur Windows 11 ; round-trip 100 \% ; carte v0 signée ;
exécutable $<$ 4 Kio ; R1–R3, R8–R9 dès le premier octet & \textbf{franchi} 26/10/2026 (D26) \\
M1 — blocs réutilisables & oracle 100 \% sur la mini-libc ; première cristallisation ;
adresses stables démontrées ; première mesure effective des budgets tokens & \textbf{franchi} 26/10/2026, 0.2.0 (D32) \\
M2 — patch à chaud & remplacement du bloc d'écho en exécution, sans interruption ;
rollback démontré ; régression post-patch nulle & \textbf{franchi} 27/10/2026, 0.3.0 (D43) \\
M3 — entrées non fiables & un million de requêtes fuzzées sans crash ; proxy
transparent ; contrats renforcés & \textbf{franchi} 28/10/2026, 0.4.0 (D74) \\
M4 — tunnel chiffré & HTTPS en transit ; échec propre sur certificat invalide & à venir \\
M5 — outil complet & interception d'un flux navigateur en VM isolée ; patch à chaud du
journal en production ; empreintes par build actives & à venir \\
M6 — portabilité inter-OS & suite de conformité verte ; \textbf{diff nul sur la couche
I} ; parité inter-profils vérifiée ; même oracle vert sur les deux profils & à venir \\
M7 — portabilité inter-ISA & round-trip 100 \% sur la nouvelle ISA ; suite de conformité
verte ; patch à chaud conforme (scratch + \code{fence.i}) ; verrou multi-profils ; export
multi-profils vérifié & à venir \\
\bottomrule
\end{tabularx}
\end{table}

M0 est le jalon existentiel — il prouve la chaîne notation, octets, exécutable, exécution
sans compilateur. M2 est le jalon de vérité — le patch à chaud a motivé l'architecture.
M6 et M7 transforment la portabilité de promesse en propriété démontrée, en isolant chaque
variable ; la suite de conformité par profil en est désormais le critère d'acceptation
explicite.

\section{Correspondance des versions}\label{sec:correspondance}

\subsection{Traçabilité des évolutions}\label{sec:evolutions}

\begin{table}[htbp]
\centering\small
\caption{Les dix-sept évolutions de la v0.4.0, tracées depuis la v0.3.}
\label{tab:correspondance}
\begin{tabularx}{\textwidth}{@{}l Y Y l@{}}
\toprule
\textbf{Évolution} & \textbf{En v0.3} & \textbf{En v0.4.0} & \textbf{Statut} \\
\midrule
Couche I normalisée & « contrats légers » sans gabarit & gabarit complet
(§\ref{sec:coucheI}) & renforcé (D19) \\
Taxonomie des défauts & absente & classes T1–T6 (§\ref{sec:taxonomie}) & nouveau (D19) \\
Threat model & 7 menaces & 13, dont 6 nouvelles & étendu (D20) \\
R10 & checksums au lancement & + signature, anti-rollback, conformité, fail-closed
explicite & renforcé (D20) \\
Optimisation prématurée & interdite implicitement & interdiction normative R11 &
explicité (D21) \\
Journal d'optimisation & pratique de démonstration & format OPT-NN & officialisé (D21) \\
Fragmentation & 4 fragments & + norme « aucune lecture totale » & consolidé (D24) \\
Localisation par intention & étape 2 implicite & balayage interdit, formalisé &
renforcé (D24) \\
Budgets tokens & 64 / 3 000 / 6 000 & localiser amendé après mesure M1, instrumentés & consolidé (D24, D31) \\
Suite de conformité profil & implicite & suite obligatoire (§\ref{sec:conformite}) &
nouveau (D22) \\
Parité inter-profils & critères dispersés & définition formelle (§\ref{sec:parite}) &
formalisé (D22) \\
Reverse-specs & absente (pratique démo) & officialisées (§\ref{sec:reversespecs}) &
officialisé (D23) \\
Pré-versions & absente (pratique \code{-dev.1}) & dev / rc / final & officialisé (D23) \\
Fixtures partagées & absente (pratique installateur) & contractuelles & officialisé
(D23) \\
Traçabilité I–C–B–tests & implicite & table normative + audit & formalisé (D24) \\
Journal des décisions & pratique & format normalisé & normalisé (D24) \\
Tableau de bord & indicateurs dispersés & métriques unifiées & nouveau (D24) \\
\bottomrule
\end{tabularx}
\end{table}

\subsection{Modifications explicitement rejetées}\label{sec:rejets}

Conformément au cahier de révision, les modifications suivantes ont été examinées et
\textbf{rejetées, donc non appliquées} :

\begin{table}[htbp]
\centering\small
\caption{Modifications examinées et rejetées.}
\label{tab:rejets}
\begin{tabularx}{\textwidth}{@{}Y Y@{}}
\toprule
\textbf{Modification} & \textbf{Motif du rejet} \\
\midrule
Modifier la hiérarchie Sécurité $>$ Performance $>$ Optimisation $>$ Coût en tokens &
invariant de révision (D19) — condition même de l'arbitrage \\
Introduire une couche d'abstraction à l'exécution (bytecode, VM) & trahit la fonction
objectif Performance ; la parité s'obtient par vérification, pas par couche \\
Promettre la portabilité universelle immédiate & la portabilité se vérifie par profil
(M6/M7), elle ne se déclare pas \\
Introduire des mécanismes de furtivité face aux EDR & hors périmètre assumé —
transparence totale de l'outil \\
Affaiblir R5, R8, R9 ou R10 & la révision ne contient que des renforcements ; R1–R9
repris à l'identique \\
Anti-rollback par contrôle en ligne au chargement & R10 doit échouer proprement hors
ligne — aucune dépendance réseau \\
Optimisation validée sans mesure & R11 : hot path + mesure avant/après + contrat stable \\
\bottomrule
\end{tabularx}
\end{table}


\subsection{Correspondance v0.4.0 $\to$ v0.5.0}\label{sec:evolutions050}

\begin{table}[htbp]
\centering\small
\caption{Les huit évolutions de la v0.5.0, tracées depuis la v0.4.0.}
\label{tab:correspondance050}
\begin{tabularx}{\textwidth}{@{}l Y Y l@{}}
\toprule
\textbf{Évolution} & \textbf{En v0.4.0} & \textbf{En v0.5.0} & \textbf{Statut} \\
\midrule
Validation terrain & constats dispersés au journal des décisions & couche normée :
boucle d'évidence, journal d'exécution, sondes contractualisées, diff chirurgical
(§\ref{sec:terrain}) & nouveau (D43) \\
Taxonomie des défauts & classes T1--T6 & + classes X1--X8, invisibles à l'émulation
(§\ref{sec:taxonomieX}) & étendu (D35--D42) \\
État des jalons & M0 planifié & M0, M1, M2 franchis — preuves terrain archivées
(§\ref{sec:jalons}) & constaté (D26, D32, D43) \\
Budgets de tokens & 40 / 3 000 / 6 000 (cibles) & 64 / 3 000 / 6 000 — corrigé et mesuré :
51 / 719 / 1 088 & corrigé (D31) \\
Suite de conformité & définie (7 éléments) & exécutée à chaque scellage rc/final — 124/124
sur le profil primaire & mesuré (D39--D43) \\
Volume d'oracle & 12 000 cas par bloc critique (référence) & campagnes effectives
consignées : 76 010 (M1), 50 533 (M2), 86 (démo) & consolidé \\
Preuves terrain & non normées & archives brutes (journaux octet-exacts, empreintes)
référencées au journal des décisions & officialisé (D43) \\
Fuzzing & mentionné (menace « fuzzing boîte noire ») & jalon M3 ouvert : fuzzing à
propriétés sur la surface d'entrées non fiables & structuré (D45) \\
\bottomrule
\end{tabularx}
\end{table}

\paragraph{Modifications explicitement rejetées (v0.5.0).} Ajouter une règle R12
« validation terrain obligatoire » — rejetée : la boucle d'évidence découle de R5 et des
critères de jalon, une règle de plus diluerait le corpus sans rien exiger de nouveau.
Étendre le fuzzing à tous les blocs scellés — rejetée : M3 cible la surface d'entrées non
fiables, les blocs à appelants de confiance relèvent des vecteurs existants. Reconstruire
les artefacts scellés 0.2.0 / 0.3.0 pour la documentation — rejetée : les balises finales
sont immuables, la documentation cite les empreintes. Fusionner les classes X dans les
classes T — rejetée : les deux familles appellent des remèdes distincts, écrire le contrat
(T) contre enrichir le modèle d'environnement et falsifier (X).

\subsection{Correspondance v0.5.0 $\to$ v0.5.1}\label{sec:evolutions051}

\begin{table}[htbp]
\centering\small
\caption{Les évolutions de la v0.5.1, tracées depuis la v0.5.0.}
\label{tab:correspondance051}
\begin{tabularx}{\textwidth}{@{}l Y Y l@{}}
\toprule
\textbf{Évolution} & \textbf{En v0.5.0} & \textbf{En v0.5.1} & \textbf{Statut} \\
\midrule
Phrases maîtresses & instruction de cérémonie seule (D27/D30) — quatre
réinitialisations du poste ont montré la limite & coffre-fort TRIADE-CF1 : conteneur
chiffré par projet, export/import, restauration post-réinitialisation & nouveau (D75) \\
État des jalons & M0–M2 constatés, M3 ouvert & M3 franchi 28/10/2026 — campagne
million (D72), conformité permanente (D73), scellage 0.4.0 (D74) & constaté (D72–D74) \\
Modèle de menace & 13 menaces, dont le vol de l'export & + vol du coffre-fort des
phrases — sans la phrase de contrôle, indéchiffrable & étendu (D75) \\
\bottomrule
\end{tabularx}
\end{table}

\paragraph{Modifications explicitement rejetées (v0.5.1).} Conserver la seule
instruction de coffre développeur — rejeté : le statu quo que les quatre réinitialisations
ont démenti. Stocker les phrases en clair dans \code{.keys/} — rejeté : l'état
opérationnel du poste n'est pas un coffre. Embarquer la graine Ed25519 dans le coffre —
rejeté : contraire à D30, jamais confiée. Confier le conteneur au seul gestionnaire de
mots de passe du développeur sans copie de coffre projet — rejeté : la continuité exige
les deux copies, le conteneur étant indéchiffrable sans la phrase.

\section{Journal des décisions}\label{sec:decisions}

Le journal complet vit dans \code{docs/DECISIONS.md} du dépôt du projet. Synthèse :
D1–D10 (v0.2) fixent l'oracle éphémère, la notation, les contrats légers, la
cryptographie, le verrou croisé, les variantes par build, l'interface système et la
granularité ; D11–D13 (infrastructure) la sauvegarde triple et systématique ; D14–D16
(v0.3) l'universalité par profils, la cryptographie paramétrable et l'export de
continuité ; D17 la régénération systématique du PDF formel ; D18 l'édition académique
(présentation IMRAD du noyau v0.3, sans valeur normative) ; D19–D24 (v0.4.0, format
normalisé) la présente révision :

\begin{itemize}[label={}, leftmargin=1.2em]
\item \textbf{D19} : couche I normalisée, taxonomie T1–T6, hiérarchie de la fonction
objectif verrouillée comme invariant de révision ;
\item \textbf{D20} : menace étendue (binaire altéré, carte corrompue, rollback,
agent défaillant, données embarquées, environnement) et R10 renforcée (checksums,
signature, anti-rollback, fail-closed explicite, vérification locale) ;
\item \textbf{D21} : R11 optimisation mesurée et journal d'optimisation OPT-NN ;
\item \textbf{D22} : suite de conformité par profil et parité inter-profils à contrat
constant ;
\item \textbf{D23} : officialisation des reverse-specs, des pré-versions dev / rc /
final et des fixtures partagées ; numérotation à trois segments ;
\item \textbf{D24} : traçabilité I $\leftrightarrow$ C $\leftrightarrow$ B
$\leftrightarrow$ tests, journal des décisions normalisé, tableau de bord de conformité.
\end{itemize}

Deux décisions postérieures à la révision complètent le tableau, sans valeur
normative : D25 confie la production du présent PDF formel à la chaîne
LaTeX/Tectonic (générateur ReportLab retiré, archives des versions antérieures
conservées) ; D26 consigne le jalon M0, première étape de la feuille de route,
menée à terme le 26 octobre 2026 (chaîne artefact implémentée, \code{hello.exe}
exécuté sur Windows 11, carte v0 signée, quatre critères de sortie acquis).

Les décisions du jalon M1 poursuivent le journal : D27 (continuité des clés après
réinitialisation du poste : régénération et centralisation, preuve par artefact
byte-identique), D28 (plan d'adressage par zones et carte projet multi-artefacts),
D29 (substrat d'émulation Unicorn pour l'oracle différentiel local), D30 (troisième
réinitialisation : D27 réappliquée, premier export de continuité généré) et D31
(budget de localisation amendé à 64 tokens : la mesure M1 a constaté que le gabarit
canonique de la présente révision mesurait déjà 43 tokens en cl100k\_base —
l'instrumentation a fait son travail de révélateur). D32 clôt le jalon : la
validation terrain de \code{selftest.exe} sur Windows 11 est constatée (ligne unique
du contrat libc-08 rendue à l'identique, 21 vecteurs sur 21), la balise devient
\code{0.2.0} finale — immuable — et le contrôle anti-rollback de la reconstruction
est étendu aux balises finales ; le jalon de vérité M2 (patch à chaud) peut
commencer.

La campagne M2 alimente ensuite le journal au rythme du terrain : D33 (base préférée
par ancre RIP-relative — image indépendante de son adresse de chargement), D34
(échelle d’auto-hébergement : la vérification migre vers le natif, chantier m7), D35
(sémantique du chargeur — FirstThunk : la première classe X de la taxonomie du
§\ref{sec:taxonomieX}), D36 à D38 (les trois premières itérations de la démonstration
graphique : lecture de structure, canal de vérité, ABI caller-saved — classes X2 à X4),
D39 à D42 (la cascade : protections mémoire, écart émulation/terrain, zone d’accueil
ABI, état hérité — classes X5 à X8), D43 (M2 clos : boucle § 5 démontrée sur machine
réelle, retrait des sondes, scellage final 0.3.0 — la validation terrain devient
couche normée). D44 consigne la présente révision v0.5.0 ; D45 ouvre le jalon M3
(contrats renforcés sur entrées non fiables, fuzzing à propriétés).

La campagne M3 poursuit le journal jusqu'à la clôture du jalon : D46 (couche I écrite
avant le code — treize spécifications du domaine http), D47 (proxy.exe scellé),
D48–D50 (lignée terrain : le superviseur un-exe et sa chaîne de lecture), D51–D71 (la
chasse aux tueurs de runs en seize générations — ABI shadow space, pile du proxy,
réponses Google, tampon de composition — jusqu'au T3 démontré de bout en bout sur
machine réelle), D72 (campagne million : 10\^6 requêtes fuzzées sans crash, reproductible),
D73 (le fuzzing rejoint la suite de conformité — sentinelle perpétuelle), D74 (M3 clos :
cristallisation 38/38 $\times$ 10/10, scellage de la balise finale 0.4.0, lignée de clés
régénérée avec preuve byte-identique). D75 consigne la présente révision v0.5.1 : le
coffre-fort des phrases maîtresses — conteneur chiffré TRIADE-CF1 centralisant toutes
les phrases par projet, sur la pile cryptographique de la carte au repos, avec autotest
de falsification et restauration post-réinitialisation.

Une décision ne se modifie jamais en silence : la révision remplace l'entrée par une note
d'invalidation explicite référençant la décision qui la supplante. La gouvernance
s'applique à elle-même — la présente révision en est la démonstration.
   % 12 Gouvernance · 13 Menaces et limites · 14 Jalons · 15 Correspondance · 16 Décisions

\end{document}
